\chapter{Representations and Characters}

\begin{center}\itshape
The group $\mathfrak S_4$ has $24$ elements, and $24=1+1+4+9+9$.\\
The group $A_5$ has $60$ elements, and $60=1+9+9+16+25$.\\
Why should the order of a group be a sum of squares?
\end{center}

A group is abstract; a group of matrices can be computed with. A \emph{representation} of a group $G$ is a way of letting $G$ act by linear maps on a vector space, as the rotations of the cube act on $\R^3$. Around 1896, answering a question of Dedekind, Frobenius discovered that the traces of these matrices, the \emph{characters}, carry a surprising amount of information, and that for each finite group there is a short list of ``irreducible'' representations, from which all others are built, whose dimensions $d_i$ satisfy $\sum d_i^2=|G|$. The proof is linear algebra, a lemma of Schur and an averaging trick.

For an abelian group the irreducible representations are one-dimensional: they are the characters of harmonic analysis. For $\Z/n\Z$ they give the discrete Fourier transform, the basis of signal processing and of the analysis of random walks (Tome~3); for $(\Z/q\Z)^*$ they are the Dirichlet characters, with which Dirichlet proved in 1837 that every progression $a,a+q,a+2q,\dots$ with $\gcd(a,q)=1$ contains infinitely many primes (Tome~5).

Throughout, $G$ is a finite group and vector spaces are finite-dimensional over $\C$.

\section{Representations}

\begin{definition}
A \emph{representation} of $G$ is a homomorphism $\rho:G\to GL(V)$, where $V$ is a complex vector space of dimension $d\ge1$, the \emph{dimension} (or degree) of $\rho$; we also say that $V$ is a representation, and write $gv$ for $\rho(g)v$. A subspace $W$ is \emph{stable} if $gW\subseteq W$ for all $g$; it is then a \emph{subrepresentation}. $V$ is \emph{irreducible} if its only stable subspaces are $0$ and $V$. A linear map $f:V\to V'$ between representations is an \emph{intertwiner} if $f(gv)=gf(v)$ for all $g,v$; an invertible one is an \emph{isomorphism}.
\end{definition}

\begin{example}\label{ex:reps}
\begin{enumerate}[label=(\arabic*)]
\item The \emph{trivial} representation: $V=\C$, $\rho(g)=1$. The \emph{sign} representation of $\mathfrak S_n$: $\rho(\sigma)=\varepsilon(\sigma)\in\{\pm1\}=GL_1(\C)\cap\R$.
\item Every one-dimensional representation is a homomorphism $\chi:G\to\C^*$, and is irreducible. For $G=\Z/n\Z$, $\chi_j(k)=e^{2i\pi jk/n}$, $0\le j<n$.
\item The \emph{permutation representation} of $\mathfrak S_n$ on $\C^n$: $\sigma e_i=e_{\sigma(i)}$. The line $\C(e_1+\dots+e_n)$ and the hyperplane $\{\sum x_i=0\}$ are stable.
\item More generally, if $G$ acts on a finite set $X$ (Chapter~3), $G$ acts on the space $\C^X$ with basis $(e_x)_{x\in X}$ by $ge_x=e_{gx}$. With $X=G$ acting on itself by left multiplication, this is the \emph{regular representation} $\C[G]$, of dimension $|G|$.
\item The symmetries of an equilateral triangle centred at $0$ give a homomorphism $\mathfrak S_3\cong D_3\to GL_2(\R)\subseteq GL_2(\C)$. It is irreducible: a stable line would be a common eigenline of the rotation by $120^\circ$, which has no real eigenvector, and of a reflection; over $\C$ the rotation's eigenlines are spanned by $(1,\mp i)$, and the reflection in the $x$-axis swaps them.
\end{enumerate}
\end{example}

\begin{think}
$\mathfrak S_3$ has order $6$, and the trivial and sign representations have dimension $1$. If $6$ is a sum of the squares of the dimensions of all irreducible representations, what is missing? Which representation above has that dimension?
\end{think}

\begin{theorem}[Maschke]\label{thm:maschke}
Every representation $V$ of $G$ has a $G$-invariant Hermitian inner product: $\langle gv,gw\rangle=\langle v,w\rangle$. Consequently every stable subspace $W$ has a stable complement $W^\perp$, and $V$ is a direct sum of irreducible subrepresentations.
\end{theorem}
\begin{proof}
Take any Hermitian inner product $(\cdot,\cdot)$ on $V$ and average it: $\langle v,w\rangle=\frac1{|G|}\sum_{h\in G}(hv,hw)$. It is Hermitian and positive definite (a sum of such), and $\langle gv,gw\rangle=\frac1{|G|}\sum_h(hgv,hgw)=\langle v,w\rangle$, because $h\mapsto hg$ is a bijection of $G$. If $W$ is stable and $v\in W^\perp$, then for $w\in W$, $\langle gv,w\rangle=\langle v,g^{-1}w\rangle=0$, so $gv\in W^\perp$. Finally, by induction on $\dim V$: if $V$ is not irreducible, $V=W\oplus W^\perp$ with both stable and of smaller dimension.
\end{proof}

So each $\rho(g)$ is unitary for $\langle\cdot,\cdot\rangle$. Since $g^n=e$ for $n=|G|$, $\rho(g)$ is a unitary matrix with $\rho(g)^n=I$: it is diagonalisable in an orthonormal basis (spectral theorem, Tome~0), with eigenvalues that are $n$-th roots of unity.

\begin{lemma}[Schur]\label{lem:schur}
Let $V,V'$ be irreducible and $f:V\to V'$ an intertwiner. Then $f=0$ or $f$ is an isomorphism. If $V=V'$, then $f=\lambda\,\mathrm{id}$ for some $\lambda\in\C$.
\end{lemma}
\begin{proof}
$\ker f$ is stable ($f(v)=0\Rightarrow f(gv)=gf(v)=0$), so it is $0$ or $V$; $f(V)$ is stable, so it is $0$ or $V'$. If $f\ne0$, then $\ker f=0$ and $f(V)=V'$. If $V=V'$, let $\lambda$ be an eigenvalue of $f$ (we are over $\C$). Then $f-\lambda\,\mathrm{id}$ is an intertwiner that is not injective, hence $0$.
\end{proof}

\begin{corollary}\label{cor:abelianirr}
If $G$ is abelian, every irreducible representation has dimension $1$.
\end{corollary}
\begin{proof}
For $h\in G$, $\rho(h)$ commutes with all $\rho(g)$, so it is an intertwiner $V\to V$, hence a scalar. Then every line is stable, and $V$, irreducible, is a line.
\end{proof}

\section{Characters and orthogonality}

\begin{definition}
The \emph{character} of a representation $\rho$ is $\chi_\rho(g)=\operatorname{tr}\rho(g)$. A \emph{class function} is a function $G\to\C$ constant on conjugacy classes; they form a vector space of dimension the number of classes, with the inner product
\[
\langle\varphi,\psi\rangle=\frac1{|G|}\sum_{g\in G}\varphi(g)\overline{\psi(g)}.
\]
\end{definition}

\begin{proposition}\label{prop:charbasic}
\begin{enumerate}[label=(\alph*)]
\item $\chi(e)=\dim V$; $\chi$ is a class function; isomorphic representations have the same character.
\item $\chi(g^{-1})=\overline{\chi(g)}$, and $|\chi(g)|\le\chi(e)$.
\item $\chi_{V\oplus W}=\chi_V+\chi_W$.
\item For a permutation representation $\C^X$, $\chi(g)$ is the number of fixed points of $g$ on $X$. In particular the character of $\C[G]$ is $|G|$ at $e$ and $0$ elsewhere.
\end{enumerate}
\end{proposition}
\begin{proof}
(a) $\operatorname{tr}I=\dim V$, and $\operatorname{tr}(ABA^{-1})=\operatorname{tr}B$ applies both to $\rho(hgh^{-1})=\rho(h)\rho(g)\rho(h)^{-1}$ and to $\rho'(g)=f\rho(g)f^{-1}$. (b) The eigenvalues $\lambda_j$ of $\rho(g)$ are roots of unity, so $\lambda_j^{-1}=\overline{\lambda_j}$; those of $\rho(g^{-1})$ are the $\lambda_j^{-1}$; and $|\sum\lambda_j|\le\sum|\lambda_j|=\dim V$. (c) In a basis adapted to $V\oplus W$ the matrices are block diagonal. (d) $\rho(g)$ is a permutation matrix, whose diagonal entry at $e_x$ is $1$ iff $gx=x$. Left multiplication by $g\ne e$ has no fixed point.
\end{proof}

\begin{theorem}[orthogonality of characters]\label{thm:orth}
Let $\chi,\chi'$ be the characters of irreducible representations $V,V'$. Then $\langle\chi,\chi'\rangle=1$ if $V\cong V'$, and $0$ otherwise.
\end{theorem}
\begin{proof}
For any linear map $u:V\to V'$, the average
\[
\tilde u=\frac1{|G|}\sum_{g\in G}\rho'(g)\,u\,\rho(g)^{-1}
\]
is an intertwiner:
\[
\rho'(h)\,\tilde u\,\rho(h)^{-1}=\frac1{|G|}\sum_g\rho'(hg)\,u\,\rho(hg)^{-1}=\tilde u.
\]
 Choose bases, write matrices with entries $\rho(g)_{ij}$, and take for $u$ the matrix unit $E_{jk}$ (entry $1$ at $(j,k)$). The $(i,l)$ entry of $\tilde u$ is $\frac1{|G|}\sum_g\rho'(g)_{ij}\rho(g^{-1})_{kl}$.

If $V\not\cong V'$, Schur gives $\tilde u=0$ for every $u$, so all these sums vanish; put $i=j$, $k=l$ and sum over $j$ and $k$: $\frac1{|G|}\sum_g\chi'(g)\chi(g^{-1})=0$, i.e.\ $\langle\chi',\chi\rangle=0$ by Proposition~\ref{prop:charbasic}(b).

If $V\cong V'$, the characters are equal, so we may take $V'=V$, $\rho'=\rho$. Schur gives $\tilde u=\lambda I$, and taking traces, $\operatorname{tr}\tilde u=\operatorname{tr}u$, so $\lambda=\operatorname{tr}(u)/d$. For $u=E_{jk}$, $\operatorname{tr}u=\delta_{jk}$, hence $\frac1{|G|}\sum_g\rho(g)_{ij}\rho(g^{-1})_{kl}=\delta_{il}\delta_{jk}/d$. Put $i=j$, $k=l$ and sum: $\frac1{|G|}\sum_g\chi(g)\chi(g^{-1})=\sum_{j,k}\delta_{jk}/d=d/d=1$.
\end{proof}

\begin{corollary}\label{cor:mult}
Let $V_1,\dots,V_r$ be pairwise non-isomorphic irreducible representations with characters $\chi_1,\dots,\chi_r$, and $V$ any representation.
\begin{enumerate}[label=(\alph*)]
\item In any decomposition of $V$ into irreducibles, the number of summands isomorphic to $V_i$ is $\langle\chi_V,\chi_i\rangle$. So two representations with the same character are isomorphic.
\item $V$ is irreducible iff $\langle\chi_V,\chi_V\rangle=1$.
\item (Regular representation) Every irreducible $V_i$ occurs in $\C[G]$, $\dim V_i$ times. Hence there are finitely many irreducible representations up to isomorphism, and if $V_1,\dots,V_r$ is the full list, $\sum_i(\dim V_i)^2=|G|$.
\end{enumerate}
\end{corollary}
\begin{proof}
(a) If $V\cong\bigoplus W_j$ with irreducible $W_j$, then $\chi_V=\sum\chi_{W_j}$ and, by Theorem~\ref{thm:orth}, $\langle\chi_V,\chi_i\rangle$ counts the $j$ with $W_j\cong V_i$. A character thus determines the multiplicities, hence the representation up to isomorphism (Maschke).

(b) If $V\cong\bigoplus V_i^{m_i}$ (all irreducibles that occur), $\langle\chi_V,\chi_V\rangle=\sum m_i^2$, which is $1$ iff $V$ is a single irreducible.

(c) By Proposition~\ref{prop:charbasic}(d), $\langle\chi_{\C[G]},\chi_i\rangle=\frac1{|G|}|G|\,\overline{\chi_i(e)}=\dim V_i$. Every irreducible occurs in the finite-dimensional $\C[G]$, so there are finitely many, and comparing dimensions in $\C[G]\cong\bigoplus V_i^{\dim V_i}$ gives $|G|=\sum(\dim V_i)^2$.
\end{proof}

\begin{theorem}\label{thm:classes}
The irreducible characters form an orthonormal basis of the space of class functions. In particular, the number of irreducible representations of $G$ is the number of its conjugacy classes.
\end{theorem}
\begin{proof}
They are orthonormal (Theorem~\ref{thm:orth}), hence linearly independent; it remains to show that a class function $f$ with $\langle f,\chi_i\rangle=0$ for all $i$ is $0$. For a representation $\rho$, let $T_\rho=\sum_g\overline{f(g)}\rho(g)$. As $\overline f$ is constant on classes, $\rho(h)T_\rho\rho(h)^{-1}=\sum_g\overline{f(g)}\rho(hgh^{-1})=T_\rho$, so $T_\rho$ is an intertwiner. On an irreducible $V_i$ it is a scalar $\lambda$ (Schur), with $\lambda\dim V_i=\operatorname{tr}T_\rho=\sum_g\overline{f(g)}\chi_i(g)=|G|\,\overline{\langle f,\chi_i\rangle}=0$. So $T_\rho=0$ on every irreducible, hence on every representation (Maschke), in particular on $\C[G]$. But $T_{\C[G]}e_e=\sum_g\overline{f(g)}e_g$, so all $f(g)=0$.
\end{proof}

\begin{corollary}[column orthogonality]\label{cor:columns}
Let $\chi_1,\dots,\chi_r$ be the irreducible characters, and $g,h\in G$. Then $\sum_i\chi_i(g)\overline{\chi_i(h)}$ is $|C_G(g)|$ if $g$ and $h$ are conjugate, and $0$ otherwise.
\end{corollary}
\begin{proof}
Let $C_1,\dots,C_r$ be the classes (there are $r$ of them, Theorem~\ref{thm:classes}), and $M$ the $r\times r$ matrix $M_{ij}=\sqrt{|C_j|/|G|}\,\chi_i(C_j)$. Row orthogonality says $M\overline M^{\,T}=I$. A square matrix with a right inverse is invertible with the same inverse, so $\overline M^{\,T}M=I$: $\sum_i\frac{\sqrt{|C_j||C_k|}}{|G|}\overline{\chi_i(C_j)}\chi_i(C_k)=\delta_{jk}$. With $|G|/|C_j|=|C_G(g)|$ for $g\in C_j$ (orbit--stabiliser, Theorem~\ref{thm:orbstab}), this is the claim.
\end{proof}

\section{Character tables}

The \emph{character table} of $G$ lists the values of the irreducible characters on the conjugacy classes; it is square by Theorem~\ref{thm:classes}. Three tools fill it in: one-dimensional representations, permutation representations (whose characters count fixed points), and the orthogonality relations. Two more constructions help. If $\chi$ is a character and $\lambda$ a one-dimensional character, $\lambda\chi$ is a character (of $\rho'(g)=\lambda(g)\rho(g)$ on the same space), irreducible with $\chi$ since $\langle\lambda\chi,\lambda\chi\rangle=\langle\chi,\chi\rangle$ ($|\lambda(g)|=1$). And if $N\trianglelefteq G$, every representation of $G/N$ gives one of $G$ by composing with $G\to G/N$; it is irreducible iff the original is, since the two have the same matrices.

\begin{example}[$\mathfrak S_3$]\label{ex:s3}
Three classes (identity, $3$ transpositions, $2$ three-cycles), so three irreducibles, with $\sum d_i^2=6$: dimensions $1,1,2$. The permutation character of $\C^3$ is $(3,1,0)$ (fixed points); minus the trivial character it gives $\chi_2=(2,0,-1)$, and $\langle\chi_2,\chi_2\rangle=\frac{4+0+2\cdot1}6=1$, so it is irreducible: it is the triangle representation of Example~\ref{ex:reps}(5).
\end{example}

\begin{example}[$\mathfrak S_4$]\label{ex:s4}
Five classes, of sizes $1,6,3,8,6$ for $e$, $(12)$, $(12)(34)$, $(123)$, $(1234)$ (Chapter~4). The permutation character of $\C^4$ is $(4,2,0,1,0)$; minus the trivial one, $\chi_3=(3,1,-1,0,-1)$, of norm $\frac{9+6+3+0+6}{24}=1$. Then $\chi_3'=\varepsilon\chi_3=(3,-1,-1,0,1)$ is another irreducible character. The quotient $\mathfrak S_4/V_4\cong\mathfrak S_3$ (Chapter~4) turns $\chi_2$ of $\mathfrak S_3$ into $(2,0,2,-1,0)$: the double transpositions lie in $V_4$, and the image of a $4$-cycle in $\mathfrak S_3$ is a transposition, since its square lies in $V_4$ but it does not. With the trivial and sign characters, these are five distinct irreducible characters, and $1+1+4+9+9=24$: the table is complete (Figure~\ref{fig:tables}). The two $3$-dimensional ones are the rotations of the cube ($\chi_3'$, through the isomorphism of Chapter~1) and the symmetries of the regular tetrahedron with vertices $e_1,\dots,e_4$ ($\chi_3$).
\end{example}

\begin{figure}[ht]
\centering
\fig{ch11-tables}
\caption{The character tables of $\mathfrak S_3$, $\mathfrak S_4$ and $A_5$. Columns are the conjugacy classes (a representative, and the size of the class in grey), rows the irreducible characters; each cell is coloured by $\chi(g)/\chi(e)$, from blue ($-1$) through white ($0$) to orange ($1$). In the table of $A_5$, $\varphi=\frac{1+\sqrt5}2$ is the golden ratio. Rows are orthonormal for the weighted inner product, and columns are orthogonal (Corollary~\ref{cor:columns}).}
\label{fig:tables}
\end{figure}

\begin{remark}
Characters prove theorems about groups that have nothing to do with linear algebra. Burnside (1904) showed with them that a group of order $p^aq^b$ ($p,q$ primes) is solvable, and Frobenius that in a group acting transitively on a set where no non-identity element fixes two points, the identity together with the elements fixing no point form a normal subgroup. No proof of the latter avoiding characters is known. The dimensions of the irreducible representations always divide $|G|$ (Exercise~\ref{ex:divides}).
\end{remark}

\section{Abelian groups: Fourier analysis}\label{sec:fourierab}

Let $A$ be a finite abelian group. By Corollary~\ref{cor:abelianirr} and Theorem~\ref{thm:classes} (each element is its own class), $A$ has exactly $|A|$ irreducible characters, all one-dimensional: the homomorphisms $\chi:A\to\C^*$. They form a group $\widehat A$ under pointwise multiplication (a product of homomorphisms to the abelian group $\C^*$ is one), the \emph{dual group}, and $|\widehat A|=|A|$. For $A=\Z/n\Z$, a homomorphism is determined by $\chi(1)$, an $n$-th root of unity, so $\widehat A=\{\chi_j:k\mapsto e^{2i\pi jk/n}\}\cong\Z/n\Z$.

\begin{theorem}[Fourier analysis on a finite abelian group]\label{thm:fourier}
For $f:A\to\C$ define $\hat f(\chi)=\langle f,\chi\rangle=\frac1{|A|}\sum_{a}f(a)\overline{\chi(a)}$. Then
\[
\begin{gathered}f=\sum_{\chi\in\widehat A}\hat f(\chi)\,\chi,\qq\frac1{|A|}\sum_{a\in A}|f(a)|^2=\sum_{\chi\in\widehat A}|\hat f(\chi)|^2,\end{gathered}
\]
and for $a\in A$,
\[
\frac1{|A|}\sum_{\chi\in\widehat A}\chi(a)=\begin{cases}1&\text{if }a=0,\\0&\text{otherwise.}\end{cases}
\]
\end{theorem}
\begin{proof}
All functions are class functions, and the characters form an orthonormal basis of them (Theorem~\ref{thm:classes}); the first formula is the expansion in an orthonormal basis, the second is Pythagoras (Parseval). The third is Corollary~\ref{cor:columns} with $h=0$: the centraliser of $a$ is $A$, and $a$ is conjugate to $0$ only if $a=0$.
\end{proof}

For $A=\Z/n\Z$, $\hat f(\chi_j)=\frac1n\sum_kf(k)e^{-2i\pi jk/n}$ is the \emph{discrete Fourier transform}. It turns convolution into multiplication (Exercise~\ref{ex:circulant}), and the fast Fourier transform of Cooley and Tukey computes it in $O(n\log n)$ operations instead of $O(n^2)$, by splitting $\Z/n\Z$ along a subgroup when $n=2^k$.

\subsection*{Dirichlet characters}

Let $q\ge1$. A \emph{Dirichlet character} modulo $q$ is a function $\chi:\Z\to\C$ obtained from a character $\chi$ of $(\Z/q\Z)^*$ by $\chi(n)=\chi(n\bmod q)$ if $\gcd(n,q)=1$, and $\chi(n)=0$ otherwise. It is periodic of period $q$ and \emph{completely multiplicative}: $\chi(mn)=\chi(m)\chi(n)$ for all $m,n$ (if $m$ or $n$ is not prime to $q$, both sides are $0$). There are $\varphi(q)$ of them; the \emph{principal} character $\chi_0$ is $1$ on integers prime to $q$.

\begin{proposition}\label{prop:dirichletorth}
Let $a$ be prime to $q$. For every integer $n$,
\[
\frac1{\varphi(q)}\sum_{\chi\bmod q}\overline{\chi(a)}\,\chi(n)=\begin{cases}1&\text{if }n\equiv a\ (q),\\0&\text{otherwise.}\end{cases}
\]
\end{proposition}
\begin{proof}
If $\gcd(n,q)>1$, every term is $0$, and $n\not\equiv a$. Otherwise $\overline{\chi(a)}\chi(n)=\chi(a^{-1}n)$ in $A=(\Z/q\Z)^*$ ($|\chi(a)|=1$, so $\overline{\chi(a)}=\chi(a)^{-1}=\chi(a^{-1})$), and Theorem~\ref{thm:fourier}, written multiplicatively, gives $1$ if $a^{-1}n=1$ in $A$ and $0$ otherwise.
\end{proof}

This identity is a \emph{filter}: to study the primes $p\equiv a\pmod q$, sum $\frac1{\varphi(q)}\overline{\chi(a)}\chi(p)$ over all primes and all $\chi$, then treat each $\chi$ separately, which is possible because $\chi$ is multiplicative and has an Euler product. This is the plan of Dirichlet's theorem, carried out in Tome~5. For $q=p$ an odd prime, $(\Z/p\Z)^*$ is cyclic of even order, so it has exactly one character of order $2$, which must be the Legendre symbol $\bigl(\frac np\bigr)$ of Tome~1: Euler's criterion $\bigl(\frac np\bigr)\equiv n^{(p-1)/2}$ shows that it is multiplicative.

\begin{example}\label{ex:dir5}
Modulo $5$, $2$ generates $(\Z/5\Z)^*=\{1,2,4,3\}=\{2^0,2^1,2^2,2^3\}$, and the four characters are determined by $\chi(2)\in\{1,i,-1,-i\}$:
\begin{center}
\small
\begin{tabular}{c|cccc}
$n\bmod5$ & $1$ & $2$ & $3$ & $4$\\\hline
$\chi_0$ & $1$ & $1$ & $1$ & $1$\\
$\chi_1$ & $1$ & $i$ & $-i$ & $-1$\\
$\chi_2$ & $1$ & $-1$ & $-1$ & $1$\\
$\chi_3$ & $1$ & $-i$ & $i$ & $-1$
\end{tabular}
\end{center}
$\chi_2$ is the Legendre symbol $\bigl(\frac n5\bigr)$: the squares modulo $5$ are $1$ and $4$.
\end{example}

\begin{beauty}
The same few lines of linear algebra describe the vibrations of a symmetric molecule (chemists read the character table of its symmetry group), the frequencies of a sampled signal (the Fourier transform on $\Z/n\Z$), the mixing time of card shuffles (Diaconis and Shahshahani used the characters of $\mathfrak S_n$ to show that about $\frac12n\log n$ random transpositions mix a deck of $n$ cards, Tome~3), and the distribution of primes in progressions (Dirichlet characters, Tome~5). Fourier analysis on $\R$ and on the circle, met in Tome~4, is the same theory for infinite abelian groups; for non-abelian Lie groups such as $SO(3)$, the next chapter's rotations, it becomes the theory of spherical harmonics and of spin in quantum mechanics.
\end{beauty}

\section{Exercises}

\begin{exercise}[$\star$]\label{ex:pairs}
$\mathfrak S_4$ acts on the $6$ two-element subsets of $\{1,2,3,4\}$. Decompose this permutation representation into irreducibles.
\end{exercise}

\begin{exercise}[$\star$]\label{ex:onedim}
Show that the one-dimensional representations of $G$ correspond to the homomorphisms $G/G'\to\C^*$, where $G'$ is the derived subgroup (Exercise~\ref{ex:derived}), so that there are $[G:G']$ of them. How many do $\mathfrak S_n$ ($n\ge2$), $A_4$ and $A_5$ have?
\end{exercise}

\begin{exercise}[$\star\star$]\label{ex:a4}
Find the character table of $A_4$. (Its classes have sizes $1,3,4,4$; use $A_4/V_4\cong\Z/3\Z$.)
\end{exercise}

\begin{exercise}[$\star\star$]\label{ex:centralchar}
Let $\rho$ be irreducible of dimension $d$. Show that $\rho(z)$ is a scalar for every $z$ in the centre $Z(G)$, and deduce that $d^2\le[G:Z(G)]$. What does this give for the quaternion group $Q_8$, and what is its character table?
\end{exercise}

\begin{exercise}[$\star\star$]\label{ex:circulant}
For $f,g:\Z/n\Z\to\C$, let $(f*g)(k)=\sum_jf(j)g(k-j)$.
\begin{enumerate}[label=(\alph*)]
\item Show that $\widehat{f*g}(\chi)=n\,\hat f(\chi)\hat g(\chi)$.
\item A \emph{circulant} matrix has entries $C_{kj}=c(k-j)$ for some $c:\Z/n\Z\to\C$. Show that each character $\chi_m$, as a vector $(\chi_m(k))_k$, is an eigenvector of $C$, and find the eigenvalues.
\item A walker on $\Z/n\Z$ ($n$ odd) moves $\pm1$ with probability $\frac12$ each. Find the Fourier coefficients of its distribution after $t$ steps, starting from $0$, and show that it converges to the uniform distribution.
\end{enumerate}
\end{exercise}

\begin{exercise}[$\star\star$]\label{ex:dirichlet8}
Write the table of Dirichlet characters modulo $8$ and modulo $12$. Check Proposition~\ref{prop:dirichletorth} for $q=8$, $a=3$, and show that all characters modulo $8$ are real.
\end{exercise}

\begin{exercise}[$\star\star\star$]\label{ex:a5}
Find the character table of $A_5$ (Figure~\ref{fig:tables}), using its permutation representations on $\{1,\dots,5\}$ and on the $10$ two-element subsets, and the orthogonality relations.
\end{exercise}

\begin{exercise}[$\star\star\star$]\label{ex:divides}
A complex number is an \emph{algebraic integer} if it is a root of a monic polynomial with integer coefficients.
\begin{enumerate}[label=(\alph*)]
\item Show that $x\in\C$ is an algebraic integer iff $xM\subseteq M$ for some non-zero finitely generated subgroup $M$ of $(\C,+)$. Deduce that sums and products of algebraic integers are algebraic integers, and that a rational algebraic integer is an integer.
\item Let $\chi$ be irreducible of dimension $d$, $C$ a conjugacy class and $g\in C$. Show that $\sum_{x\in C}\rho(x)$ is the scalar $\omega_C=|C|\chi(g)/d$, and that $\omega_C$ is an algebraic integer.
\item Deduce that $d$ divides $|G|$.
\end{enumerate}
\end{exercise}

\section{Hints}

\begin{hint}{ex:pairs}
Count the fixed pairs of each class, then take inner products with the five rows of the table.
\end{hint}
\begin{hint}{ex:onedim}
$\C^*$ is abelian. For $\mathfrak S_n$, $A_4$ and $A_5$, find the derived subgroups with Chapters~4 and 5.
\end{hint}
\begin{hint}{ex:a4}
Three one-dimensional characters come from $\Z/3\Z$; then $12=1+1+1+9$. For the last row use column orthogonality, or the permutation representation on $4$ points.
\end{hint}
\begin{hint}{ex:centralchar}
Schur. Then $1=\langle\chi,\chi\rangle\ge\frac1{|G|}\sum_{z\in Z(G)}|\chi(z)|^2$. For $Q_8$: $|Z|=2$, $8=1+1+1+1+4$.
\end{hint}
\begin{hint}{ex:circulant}
(a) Expand and substitute $k=j+l$. (b) $(C\chi_m)(k)=\sum_jc(k-j)\chi_m(j)$; substitute $j=k-l$. (c) The distribution after $t$ steps is a $t$-fold convolution.
\end{hint}
\begin{hint}{ex:dirichlet8}
$(\Z/8\Z)^*\cong(\Z/2\Z)^2$, generated by $3$ and $5$; $(\Z/12\Z)^*=\{1,5,7,11\}$, likewise.
\end{hint}
\begin{hint}{ex:a5}
The first gives $\psi_4$, the second $\psi_4+\psi_5+\mathbf1$. Then $60=1+16+25+d^2+d'^2$; use column orthogonality on each column, and the fact that the values on $5$-cycles are real ($g$ and $g^{-1}$ are conjugate in $A_5$).
\end{hint}
\begin{hint}{ex:divides}
(a) If $M=\Z m_1+\dots+\Z m_k$ and $xm_i=\sum a_{ij}m_j$, then $x$ is an eigenvalue of the integer matrix $(a_{ij})$. (b) The products $e_Ce_D$, $e_C=\sum_{x\in C}\rho(x)$, are integer combinations of the $e_E$. (c) $\sum_C\omega_C\overline{\chi(g_C)}=|G|/d$.
\end{hint}

\section{Solutions}

\begin{solution}{ex:pairs}
\idea Characters of permutation representations are fixed-point counts; multiplicities are inner products.
\writeup On the classes of $e$, $(12)$, $(12)(34)$, $(123)$, $(1234)$ (of sizes $1$, $6$, $3$, $8$, $6$), the numbers of fixed pairs are $6$; $2$ ($\{1,2\},\{3,4\}$); $2$ ($\{1,2\},\{3,4\}$); $0$; $0$. So $\chi=(6,2,2,0,0)$ and
\[
\begin{gathered}
\langle\chi,\mathbf1\rangle=\tfrac{6+12+6}{24}=1,\qq\langle\chi,\varepsilon\rangle=\tfrac{6-12+6}{24}=0,\\
\langle\chi,\chi_2\rangle=\tfrac{12+0+12}{24}=1,\qq\langle\chi,\chi_3\rangle=\tfrac{18+12-6}{24}=1,\\
\langle\chi,\chi_3'\rangle=\tfrac{18-12-6}{24}=0.
\end{gathered}
\]
So $\chi=\mathbf1+\chi_2+\chi_3$, and the dimensions add up: $1+2+3=6$.
\end{solution}

\begin{solution}{ex:onedim}
\idea A homomorphism to an abelian group kills all commutators.
\writeup A one-dimensional representation is a homomorphism $\lambda:G\to\C^*$. As $\C^*$ is abelian, $\lambda(xyx^{-1}y^{-1})=1$, so $\lambda$ is trivial on $G'$ and factors through $G/G'$; conversely a homomorphism $G/G'\to\C^*$ composed with $G\to G/G'$ is one. The group $G/G'$ is abelian, so it has $|G/G'|$ characters (Section~\ref{sec:fourierab}), and distinct ones give distinct $\lambda$. For $\mathfrak S_n$, $n\ge3$: $G'\subseteq A_n$ because $\mathfrak S_n/A_n$ is abelian; and $[(a\,b),(a\,c)]=\bigl((a\,b)(a\,c)\bigr)^2$ is the square of a $3$-cycle, hence a $3$-cycle, and every $3$-cycle arises this way (it is the square of its inverse). The $3$-cycles generate $A_n$ (Proposition~\ref{prop:gens}), so $G'=A_n$ and there are two one-dimensional representations, trivial and sign; the same holds for $n=2$, as $\mathfrak S_2$ is abelian of order $2$. For $A_4$: $G'\subseteq V_4$ since $A_4/V_4\cong\Z/3\Z$ is abelian, and $G'\ne\{e\}$ since $A_4$ is not abelian. A normal subgroup is a union of classes, of sizes $1,3,4,4$, so a normal subgroup of order $2$ or $4$ inside $V_4$ must have order $1+3=4$: $G'=V_4$, and there are three. For $A_5$, simple and not abelian, $G'$ is a non-trivial normal subgroup, so $G'=A_5$: only the trivial representation.
\end{solution}

\begin{solution}{ex:a4}
\idea Lift the characters of $\Z/3\Z$, then let orthogonality give the last row.
\writeup Let $\omega=e^{2i\pi/3}$. The classes are $e$, the double transpositions ($V_4\setminus\{e\}$, $3$ elements), the class of $(1\,2\,3)$ and that of $(1\,3\,2)$ ($4$ each). The quotient $A_4\to A_4/V_4\cong\Z/3\Z$ sends $(1\,2\,3)$ to a generator; the three characters of $\Z/3\Z$ give $\lambda_k=(1,1,\omega^k,\omega^{2k})$, $k=0,1,2$. The remaining irreducible has $d^2=12-3=9$, $d=3$. It is the permutation character $(4,0,1,1)$ minus the trivial one, $\chi=(3,-1,0,0)$, of norm $\frac{9+3+0+0}{12}=1$:
\begin{center}
\small
\begin{tabular}{c|cccc}
& $e$ & $(12)(34)$ & $(123)$ & $(132)$\\
size & $1$ & $3$ & $4$ & $4$\\\hline
$\lambda_0$ & $1$ & $1$ & $1$ & $1$\\
$\lambda_1$ & $1$ & $1$ & $\omega$ & $\omega^2$\\
$\lambda_2$ & $1$ & $1$ & $\omega^2$ & $\omega$\\
$\chi$ & $3$ & $-1$ & $0$ & $0$
\end{tabular}
\end{center}
The $3$-dimensional representation is the rotation group of the regular tetrahedron.
\end{solution}

\begin{solution}{ex:centralchar}
\idea A central element commutes with everything, so Schur makes it a scalar of modulus $1$.
\writeup For $z\in Z(G)$, $\rho(z)\rho(g)=\rho(zg)=\rho(gz)=\rho(g)\rho(z)$, so $\rho(z)$ is an intertwiner of the irreducible $V$, a scalar $\lambda I$, with $|\lambda|=1$ (a root of unity). So $|\chi(z)|=d$, and
\[
1=\langle\chi,\chi\rangle\ge\frac1{|G|}\sum_{z\in Z(G)}|\chi(z)|^2=\frac{|Z(G)|\,d^2}{|G|}.
\]
For $Q_8=\{\pm1,\pm i,\pm j,\pm k\}$: $Z=\{\pm1\}$, so $d^2\le4$. The five classes are
\[
\{1\},\{-1\},\{\pm i\},\{\pm j\},\{\pm k\}
\]
($i$ and $-i=jij^{-1}$ are conjugate); $8=\sum d_i^2$ with five terms $d_i\in\{1,2\}$ forces $1,1,1,1,4$. $Q_8/\{\pm1\}\cong(\Z/2\Z)^2$ gives four one-dimensional characters: $\pm1$ on the classes of $i,j,k$ with an even number of minus signs. The last one, $\chi$, has $\chi(1)=2$. Its other values follow from column orthogonality with the column of $1$ (Corollary~\ref{cor:columns}): the four one-dimensional characters are $1$ at $-1$, so $4+2\chi(-1)=0$ and $\chi(-1)=-2$; and at $i$ they sum to $1+1-1-1=0$, so $\chi(i)=0$, and likewise $\chi(j)=\chi(k)=0$. It is realised by the quaternions of Chapter~12 acting on $\C^2$. Curiously, $D_4$ has the same character table, although $D_4\not\cong Q_8$: the character table does not determine the group.
\end{solution}

\begin{solution}{ex:circulant}
\idea Characters are the common eigenvectors of all translations, so they diagonalise everything that commutes with translations.
\writeup (a) $n\,\widehat{f*g}(\chi)=\sum_k\sum_jf(j)g(k-j)\overline{\chi(k)}$; with $k=j+l$ and $\overline{\chi(j+l)}=\overline{\chi(j)}\,\overline{\chi(l)}$, this is $\sum_jf(j)\overline{\chi(j)}\sum_lg(l)\overline{\chi(l)}=n\hat f(\chi)\cdot n\hat g(\chi)$.

(b) $(C\chi_m)(k)=\sum_jc(k-j)\chi_m(j)=\sum_lc(l)\chi_m(k-l)=\Bigl(\sum_lc(l)\overline{\chi_m(l)}\Bigr)\chi_m(k)$. So the eigenvalue is $\sum_lc(l)e^{-2i\pi ml/n}=n\hat c(\chi_m)$, and the $n$ characters, orthogonal, diagonalise $C$.

(c) One step has distribution $s=\frac12(\delta_1+\delta_{-1})$, with $n\hat s(\chi_m)=\frac12(e^{-2i\pi m/n}+e^{2i\pi m/n})=\cos\frac{2\pi m}n$. After $t$ steps the distribution is $s^{*t}$, and by (a) $n\widehat{s^{*t}}(\chi_m)=\cos^t\frac{2\pi m}n$. For $m\ne0$ and $n$ odd, $|\cos\frac{2\pi m}n|<1$ ($\frac{2\pi m}n$ is not a multiple of $\pi$), so these tend to $0$, and by Fourier inversion $s^{*t}(k)=\sum_m\widehat{s^{*t}}(\chi_m)\chi_m(k)\to\frac1n$, the uniform distribution. The slowest term, $\cos^t\frac{2\pi}n\approx e^{-2\pi^2t/n^2}$, shows that about $n^2$ steps are needed (Tome~3). For even $n$, $m=n/2$ gives $(-1)^t$: the walk alternates between even and odd points and never converges.
\end{solution}

\begin{solution}{ex:dirichlet8}
\idea Decompose the group into cyclic factors and choose independently the values on generators.
\writeup Modulo $8$: $3^2\equiv5^2\equiv1$ and $15\equiv7$, so $(\Z/8\Z)^*=\{1,3,5,7\}\cong(\Z/2\Z)^2$, and a character is fixed by $\chi(3),\chi(5)\in\{\pm1\}$, with $\chi(7)=\chi(3)\chi(5)$. All values are $\pm1$: the characters are real. Modulo $12$: $5^2\equiv7^2\equiv1$, $5\cdot7\equiv11$, same structure:
\begin{center}
\small
\begin{tabular}{c|cccc}
$\bmod8$ & $1$ & $3$ & $5$ & $7$\\\hline
$\chi_0$ & $1$ & $1$ & $1$ & $1$\\
$\chi_1$ & $1$ & $-1$ & $1$ & $-1$\\
$\chi_2$ & $1$ & $1$ & $-1$ & $-1$\\
$\chi_3$ & $1$ & $-1$ & $-1$ & $1$
\end{tabular}\qquad
\begin{tabular}{c|cccc}
$\bmod12$ & $1$ & $5$ & $7$ & $11$\\\hline
$\chi_0$ & $1$ & $1$ & $1$ & $1$\\
$\chi_1$ & $1$ & $-1$ & $1$ & $-1$\\
$\chi_2$ & $1$ & $1$ & $-1$ & $-1$\\
$\chi_3$ & $1$ & $-1$ & $-1$ & $1$
\end{tabular}
\end{center}
For $a=3$: $\frac14\sum_\chi\chi(3)\chi(n)$ (real characters) at $n=1,3,5,7$ is $\frac14(1-1+1-1)=0$, $\frac14(1+1+1+1)=1$, $\frac14(1-1-1+1)=0$, $\frac14(1+1-1-1)=0$; and $0$ for even $n$. It picks out $n\equiv3\pmod8$.
\end{solution}

\begin{solution}{ex:a5}
\idea Two permutation representations give $\psi_4$ and $\psi_5$; orthogonality pins down the two missing rows, and the golden ratio appears as a root of $x^2-x-1$.
\writeup Classes (Theorem~\ref{thm:a5}): $e$, $(12)(34)$, $(123)$, $(12345)$, $(13524)$, of sizes $1,15,20,12,12$. On $\{1,\dots,5\}$ the fixed-point counts are $(5,1,2,0,0)$; minus $\mathbf1$, $\psi_4=(4,0,1,-1,-1)$, of norm $\frac{16+0+20+12+12}{60}=1$. On the $10$ pairs, fixed pairs: $10$; $2$ ($\{1,2\}$, $\{3,4\}$); $1$ ($\{4,5\}$); $0$; $0$. Minus $\mathbf1$ and $\psi_4$: $\psi_5=(5,1,-1,0,0)$, of norm $\frac{25+15+20}{60}=1$. The remaining irreducibles satisfy $\sum d^2=60-1-16-25=18$, with five classes, so there are two, of dimension $3$ ($9+9$): $\psi_3=(3,a,b,c,c')$ and $\psi_3'=(3,a',b',e,e')$.

Column orthogonality (Corollary~\ref{cor:columns}) between the column of $g$ and the column of $e$ gives $\sum_\chi\chi(e)\chi(g)=0$, and the column of $g$ has squared norm $|C_G(g)|=60/|\text{class}|$. The values are real: each $g\in A_5$ is conjugate to $g^{-1}$ (for the $5$-cycle, $(2\,5)(3\,4)$ conjugates $(12345)$ to $(15432)$), so $\chi(g)=\chi(g^{-1})=\overline{\chi(g)}$.
\begin{itemize}
\item $(12)(34)$, centraliser of order $4$: $1+0+5+3a+3a'=0$ and $1+0+1+a^2+a'^2=4$; so $a+a'=-2$, $a^2+a'^2=2$, $a=a'=-1$.
\item $(123)$, centraliser of order $3$: $1+4-5+3b+3b'=0$ and $1+1+1+b^2+b'^2=3$; so $b=b'=0$.
\item $(12345)$, centraliser of order $5$: $1-4+0+3c+3e=0$ and $1+1+0+c^2+e^2=5$; so $c+e=1$, $c^2+e^2=3$, $ce=-1$: $c,e$ are the roots of $x^2-x-1$, i.e.\ $\varphi$ and $1-\varphi$. The same holds for $c',e'$; and $\langle\psi_3,\mathbf1\rangle=0$ gives $3-15+0+12c+12c'=0$, $c+c'=1$, so $c'=1-c$.
\end{itemize}
Choosing $\psi_3(12345)=\varphi$: $\psi_3=(3,-1,0,\varphi,1-\varphi)$, $\psi_3'=(3,-1,0,1-\varphi,\varphi)$. $\psi_3$ is the rotation representation of $A_5$ as the symmetry group of the icosahedron (Chapter~12): a rotation by $\frac{2\pi}5$ has trace $1+2\cos\frac{2\pi}5=\varphi$.
\end{solution}

\begin{solution}{ex:divides}
\idea Integrality is stability of a finitely generated lattice; the central elements $e_C$ span such a lattice, and they act by scalars.
\writeup (a) If $x^n+a_{n-1}x^{n-1}+\dots+a_0=0$ with $a_i\in\Z$, then $M=\Z+\Z x+\dots+\Z x^{n-1}$ satisfies $xM\subseteq M$ and $1\in M$. Conversely, if $M=\Z m_1+\dots+\Z m_k\ne0$ and $xm_i=\sum_ja_{ij}m_j$ with $a_{ij}\in\Z$, then $(xI-A)m=0$ for the non-zero vector $m=(m_i)$, so $\det(xI-A)=0$: $x$ is a root of the monic integer polynomial $\det(XI-A)$. If $x,y$ are algebraic integers, of degrees $n,n'$ as above, $M=\sum_{i<n,j<n'}\Z x^iy^j$ is finitely generated, contains $1$, and is stable under multiplication by $x$, $y$, hence by $x+y$ and $xy$. A rational root $a/b$ (in lowest terms) of a monic integer polynomial has $b\mid1$ (Tome~0, rational roots).

(b) $e_C=\sum_{x\in C}\rho(x)$ commutes with every $\rho(g)$ (conjugation by $g$ permutes $C$), so it is a scalar $\omega_C$ (Schur); taking traces, $d\,\omega_C=|C|\chi(g)$. Summing over pairs,
\[
e_Ce_D=\sum_{x\in C,\,y\in D}\rho(xy)=\sum_Ea_{CDE}\,e_E,
\]
where $a_{CDE}\ge0$ is the number of pairs $(x,y)\in C\times D$ with $xy=z$, for a fixed $z\in E$ (the same for all $z\in E$, by conjugation). So $\omega_C\omega_D=\sum_Ea_{CDE}\,\omega_E$. So $M=\sum_E\Z\omega_E$ is finitely generated, non-zero ($\omega_{\{e\}}=1$), and $\omega_CM\subseteq M$: $\omega_C$ is an algebraic integer by (a).

(c) $\overline{\chi(g)}$ is a sum of roots of unity, an algebraic integer. Grouping by classes,
\[
\begin{aligned}\frac{|G|}d=\frac{|G|}d\langle\chi,\chi\rangle&=\frac1d\sum_C|C|\chi(g_C)\overline{\chi(g_C)}\eqbr=\sum_C\omega_C\overline{\chi(g_C)},\end{aligned}
\]
an algebraic integer by (a), and rational. So $|G|/d\in\Z$.
\end{solution}

\section*{Sources and further reading}

\begin{itemize}
\item G.~Frobenius, \emph{Über Gruppencharaktere}, Sitzungsber.\ Preuss.\ Akad.\ Wiss.\ Berlin (1896), 985--1021. The birth of character theory.
\item H.~Maschke, \emph{Beweis des Satzes, dass diejenigen endlichen linearen Substitutionsgruppen, in welchen einige durchgehends verschwindende Coefficienten auftreten, intransitiv sind}, Math.\ Ann.\ 52 (1899), 363--368. Complete reducibility.
\item I.~Schur, \emph{Neue Begründung der Theorie der Gruppencharaktere}, Sitzungsber.\ Preuss.\ Akad.\ Wiss.\ Berlin (1905), 406--432. Schur's lemma and the proofs of this chapter.
\item J.-P.~Serre, \emph{Linear Representations of Finite Groups}, Springer, 1977, Part~I. The model for this chapter, in forty pages.
\item G.~James and M.~Liebeck, \emph{Representations and Characters of Groups}, 2nd ed., Cambridge University Press, 2001. Many character tables worked out, including $A_5$.
\item P.~Diaconis, \emph{Group Representations in Probability and Statistics}, IMS Lecture Notes 11, 1988. Characters and Fourier analysis for random walks and shuffles.
\item J.~W.~Cooley and J.~W.~Tukey, \emph{An algorithm for the machine calculation of complex Fourier series}, Math.\ Comp.\ 19 (1965), 297--301. The fast Fourier transform.
\item P.~G.~L.~Dirichlet, \emph{Beweis des Satzes, dass jede unbegrenzte arithmetische Progression \dots\ unendlich viele Primzahlen enthält}, Abh.\ Königl.\ Preuss.\ Akad.\ Wiss.\ (1837), 45--81. Characters of $(\Z/q\Z)^*$, before the word existed.
\end{itemize}
